\chapter{[序列] 数据结构}
\section{序列}
\textcolor{blue}{\textbf{序列}}: 一块可存放多个值的连续内存空间，这些值按一定顺序排列，可通过每个值所在位置的编号（称为索引）访问它们。

\subsection{序列内置函数}
下面的函数适用于所有序列类型，包括字符串、列表、元组等。

\subsubsection{len 序列的长度}

\begin{enumerate}[label=\arabic*)]
    \item 字符串
\begin{lstlisting}[style=demostyle,numbers=none]
chinese_zodiac = "鼠牛虎兔龙蛇马羊猴鸡狗猪"
print(len(chinese_zodiac))
>>> 12
\end{lstlisting}

    \item 列表
\begin{lstlisting}[style=demostyle,numbers=none]
name = ['张三', '李四', '王五']
print(len(name))
>> 3
\end{lstlisting}

    \item 元组
\begin{lstlisting}[style=demostyle,numbers=none]
name = ('张三', '李四', '王五')
print(len(name))
>> 3
\end{lstlisting}

    \item 字典
\begin{lstlisting}[style=demostyle,numbers=none]
name = {'name': '张三', 'age': 18}
print(len(name))
>> 2
\end{lstlisting}

    \item 集合
\begin{lstlisting}[style=demostyle,numbers=none]
name = {'张三', '李四', '王五'}
print(len(name))
>> 3
\end{lstlisting}
\end{enumerate}


\subsubsection{max 序列的最大值}

\begin{enumerate}[label=\arabic*)]
    \item 字符串
\begin{lstlisting}[style=demostyle,numbers=none]
chinese_zodiac = "鼠牛虎兔龙蛇马羊猴鸡狗猪"
print(max(chinese_zodiac))
>>> 龙

my_string = "hello"
print(max(my_string))  # 输出: 'o'
>>> o
\end{lstlisting}

    \item 列表
\begin{lstlisting}[style=demostyle,numbers=none]
name = ['张三', '李四', '王五']
print(max(name))
>>> 王五

my_list = [1, 2, 3, 4, 5]
print(max(my_list))  # 输出: 5
>>> 5
\end{lstlisting}

    \item 元组
\begin{lstlisting}[style=demostyle,numbers=none]
name = ('张三', '李四', '王五')
print(max(name))
>>> 王五

my_tuple = (1, 2, 3)
print(max(my_tuple))  # 输出: 3
>>> 3
\end{lstlisting}

    \item 字典
\begin{lstlisting}[style=demostyle,numbers=none]
my_dict = {'a': 1, 'b': 2, 'c': 3}
print(max(my_dict))  # 输出: ('c', 3)
>>> c
print(max(my_dict.values()))
>>> 3
\end{lstlisting}

    \item 集合
\begin{lstlisting}[style=demostyle,numbers=none]
my_set = {1, 2, 3, 4}
print(max(my_set))  # 输出: 4
>>> 4
\end{lstlisting}
\end{enumerate}


\subsubsection{min 序列的最小值}
\begin{enumerate}[label=\arabic*)]
    \item 字符串
\begin{lstlisting}[style=demostyle,numbers=none]
chinese_zodiac = "鼠牛虎兔龙蛇马羊猴鸡狗猪"
print(min(chinese_zodiac))  # 输出: '鼠'
>>> 鼠
\end{lstlisting}

    \item 列表
\begin{lstlisting}[style=demostyle,numbers=none]
name = ['张三', '李四', '王五']
print(min(name))  # 输出: '张三'
>>> 张三
\end{lstlisting}

    \item 元组
\begin{lstlisting}[style=demostyle,numbers=none]
name = ('张三', '李四', '王五')
print(min(name))  # 输出: '张三'
>>> 张三
\end{lstlisting}

    \item 字典
\begin{lstlisting}[style=demostyle,numbers=none]
my_dict = {'a': 1, 'b': 2, 'c': 3}
print(min(my_dict))  # 输出: ('a', 1)
>>> a
print(min(my_dict.values()))
>>> 1
\end{lstlisting}

    \item 集合
\begin{lstlisting}[style=demostyle,numbers=none]
my_set = {1, 2, 3, 4}
print(min(my_set))  # 输出: 1
>>> 1
\end{lstlisting}

\end{enumerate}


\subsubsection{sum 返回序列的和}
\begin{enumerate}[label=\arabic*)]
    \item 字符串

    字符串不能使用 \codeinline{sum} 函数。

    \item 列表
\begin{lstlisting}[style=demostyle,numbers=none]
my_list = [1, 2, 3, 4, 5]
print(sum(my_list))  # 输出: 15
\end{lstlisting}

    \item 元组
\begin{lstlisting}[style=demostyle,numbers=none]
my_tuple = (1, 2, 3)
print(sum(my_tuple))  # 输出: 6
\end{lstlisting}

    \item 字典

    字典只能对值进行求和。
\begin{lstlisting}[style=demostyle,numbers=none]
# 字典（需要对值进行求和）
my_dict = {'a': 1, 'b': 2, 'c': 3}
print(sum(my_dict.values()))  # 输出: 6
\end{lstlisting}
    \item 集合

\begin{lstlisting}[style=demostyle,numbers=none]
my_set = {1, 2, 3, 4}
print(sum(my_set))  # 输出: 10
\end{lstlisting}
    
\end{enumerate}


\subsubsection{sorted 返回排序后的序列}
\begin{enumerate}[label=\arabic*)]
    \item 字符串
\begin{lstlisting}[style=demostyle,numbers=none]
my_string = "cba"
print(sorted(my_string))  
>>> ['a', 'b', 'c']
\end{lstlisting}

    \item 列表
\begin{lstlisting}[style=demostyle,numbers=none]
my_list = [3, 1, 4, 1, 5]
print(sorted(my_list))  
>>> [1, 1, 3, 4, 5]
\end{lstlisting}

    \item 元组（排序后转为列表）
\begin{lstlisting}[style=demostyle,numbers=none]
my_tuple = (3, 1, 4, 1, 5)
print(sorted(my_tuple))  
>>> [1, 1, 3, 4, 5]
\end{lstlisting}

    \item 字典（按键排序）
\begin{lstlisting}[style=demostyle,numbers=none]
my_dict = {'b': 2, 'a': 1, 'c': 3}
print(sorted(my_dict))
>>> ['a', 'b', 'c']
\end{lstlisting}

    \item 集合（排序后转为列表）
\begin{lstlisting}[style=demostyle,numbers=none]
my_set = {3, 1, 4}
print(sorted(my_set))
>>> [1, 3, 4]
\end{lstlisting}
\end{enumerate}

\subsubsection{reversed 返回逆序的序列}
\begin{enumerate}[label=\arabic*)]
    \item 字符串
\begin{lstlisting}[style=demostyle,numbers=none]
my_string = "hello"
print(''.join(reversed(my_string)))  # 输出: "olleh"
>>> olleh
\end{lstlisting}

    \item 列表
\begin{lstlisting}[style=demostyle,numbers=none]
my_list = [1, 2, 3, 4, 5]
print(list(reversed(my_list)))  # 输出: [5, 4, 3, 2, 1]
>>> [5, 4, 3, 2, 1]
\end{lstlisting}

    \item 元组
\begin{lstlisting}[style=demostyle,numbers=none]
my_tuple = (1, 2, 3, 4, 5)
print(tuple(reversed(my_tuple)))  # 输出: (5, 4, 3, 2, 1)
>>> (5, 4, 3, 2, 1)
\end{lstlisting}

    \item 字典（不能直接逆序，但可以逆序键或值）
\begin{lstlisting}[style=demostyle,numbers=none]
my_dict = {'a': 1, 'b': 2, 'c': 3}
print(list(reversed(my_dict)))  # 输出: ['c', 'b', 'a'] (逆序键)
>>> ['c', 'b', 'a']
\end{lstlisting}

    \item 集合

    集合是无序的，所以不能使用 \codeinline{reversed} 函数。
\end{enumerate}


\subsubsection{enumerate 返回索引和元素组成的元组}
\begin{enumerate}[label=\arabic*)]
    \item 字符串

    字符串不能使用 \codeinline{enumerate} 函数。

    \item 列表
\begin{lstlisting}[style=demostyle,numbers=none]
>>> my_list = ['a', 'b', 'c']
>>> for index, value in enumerate(my_list):
...     print(index, value)
...
0 a
1 b
2 c
\end{lstlisting}

    \item 元组
\begin{lstlisting}[style=demostyle,numbers=none]
>>> my_tuple = ('a', 'b', 'c')
>>> for index, value in enumerate(my_tuple):
...     print(index, value)
...
0 a
1 b
2 c
\end{lstlisting}

    \item 字典

    字典不能使用 \codeinline{enumerate} 函数。

    \item 集合

    集合不能使用 \codeinline{enumerate} 函数。
\end{enumerate}


\subsection{序列的通用操作}

\subsubsection{[] 索引访问}

\begin{enumerate}[label=\arabic*)]
    \item 字符串
\begin{lstlisting}[style=demostyle,numbers=none]
>>> chinese_zodiac = "鼠牛虎兔龙蛇马羊猴鸡狗猪"
>>> print(chinese_zodiac[0])
鼠
>>> print(chinese_zodiac[0:4])
鼠牛虎兔
>>> print(chinese_zodiac[-1])
猪
>>> print(chinese_zodiac[-3:-1])
鸡狗
>>> print(chinese_zodiac[-3:])
鸡狗猪
>>> print(chinese_zodiac[0:-1:2])
鼠虎龙马猴狗
\end{lstlisting}

    \item 列表
\begin{lstlisting}[style=demostyle,numbers=none]
变量名 = [值1, 值2, 值3,...]

name = ['张三', '李四', '王五']
print(name[1])
>> 李四
\end{lstlisting}

    \item 元组
\begin{lstlisting}[style=demostyle,numbers=none]
变量名 = (值1, 值2, 值3,...)

name = ('张三', '李四', '王五')
print(name[1])
>> 李四
\end{lstlisting}

    \item 字典
\begin{lstlisting}[style=demostyle,numbers=none]
变量名 = {键1:值1, 键2:值2, 键3:值3,...}

name = {'name':'张三', 'age':18, 'sex':'男'}
print(name['name'])
>> 张三
\end{lstlisting}

    \item 集合
\begin{lstlisting}[style=demostyle,numbers=none]
变量名 = {值1, 值2, 值3,...}

name = {'张三', '李四', '王五'}
print(name[1])
>> 李四
\end{lstlisting}
\end{enumerate}

\subsubsection{index 查找元素的索引}
\begin{enumerate}[label=\arabic*)]
    \item 字符串

    字符串没有 \codeinline{index} 方法，但是可以使用 \codeinline{find} 方法。

\begin{lstlisting}[style=demostyle,numbers=none]
# 字符串
my_string = "Hello, World!"
index_of_world = my_string.find("World")  # 查找子字符串 "World" 的索引
print(index_of_world)  # 输出: 7

# 未找到时返回 -1
index_of_python = my_string.find("Python")
print(index_of_python)  # 输出: -1
\end{lstlisting}

    \item 列表
\begin{lstlisting}[style=demostyle,numbers=none]
变量名 = [1, 2, 3, 4, 5]

print(变量名.index(3))
>> 2

# 查找不存在的元素会抛出 ValueError
# print(my_list.index(6))  # ValueError: 6 is not in list
\end{lstlisting}

    \item 元组
\begin{lstlisting}[style=demostyle,numbers=none]
变量名 = (1, 2, 3, 4, 5)

print(变量名.index(3))
>> 2
\end{lstlisting}

    \item 字典

    字典没有 \codeinline{index} 方法，因为字典是键值对集合，而不是简单的元素序列。通常，我们会直接通过键来访问值。

\begin{lstlisting}[style=demostyle,numbers=none]
变量名 = {'key1': 'value1', 'key2': 'value2'}

print(变量名.keys())
>> dict_keys(['key1', 'key2'])

print(变量名.values())
>> dict_values(['value1', 'value2'])

print(变量名.items())
>> dict_items([('key1', 'value1'), ('key2', 'value2')])
\end{lstlisting}

    \item 集合  

    集合没有 \codeinline{index} 方法，因为集合是无序的。
\end{enumerate}


\subsubsection{count 统计元素出现的次数}
\begin{enumerate}[label=\arabic*)]
    \item 字符串
\begin{lstlisting}[style=demostyle,numbers=none]
my_string = "Hello, world! Hello, everyone!"
count_of_hello = my_string.count("Hello")  # 统计子字符串 "Hello" 出现的次数
print(count_of_hello) 
>>> 2
\end{lstlisting}

    \item 列表
\begin{lstlisting}[style=demostyle,numbers=none]
my_list = [1, 2, 3, 1, 2, 3, 1, 2, 3]
count_of_1 = my_list.count(1)  # 统计元素 1 出现的次数
print(count_of_1)
>>> 3
\end{lstlisting}

    \item 元组
\begin{lstlisting}[style=demostyle,numbers=none]
my_tuple = (1, 2, 3, 1, 2, 3, 1, 2, 3)
count_of_1 = my_tuple.count(1)  # 统计元素 1 出现的次数
print(count_of_1)
>>> 3
\end{lstlisting}

    \item 字典

    字典没有 \codeinline{count} 方法，因为字典是键值对集合，而不是简单的元素序列。

    \item 集合

    集合没有 \codeinline{count} 方法，因为集合中的元素是唯一的，每个元素最多只出现一次。
\end{enumerate}


\subsubsection{切片}
\begin{lstlisting}[style=demostyle,numbers=none]
list[start:stop:step]
# start: 起始位置
# stop: 终止位置
# step: 步长
\end{lstlisting}

返回 \textcolor{violet}{左闭右开}区间
\begin{enumerate}[label=\arabic*)]
    \item \codeinline{[start:stop:step]}: 返回一个新的序列
    \item \codeinline{start}: 起始位置
    \item \codeinline{stop}: 终止位置
    \item \codeinline{step}: 步长
\end{enumerate}

\begin{enumerate}[label=\arabic*)]
    \item 字符串
\begin{lstlisting}[style=demostyle,numbers=none]
>>> chinese_zodiac = "鼠牛虎兔龙蛇马羊猴鸡狗猪"
>>> print(chinese_zodiac[0:4])
鼠牛虎兔
>>> print(chinese_zodiac[0:-1:2])
鼠虎龙马猴狗
\end{lstlisting}

    \item 列表
\begin{lstlisting}[style=demostyle,numbers=none]
变量名 = [值1, 值2, 值3,...]
变量名[起始位置:终止位置:步长]

name = ['张三', '李四', '王五', '乔治', '佩奇']
print(name[1:3])
>> ['李四', '王五']
print(name[3:])
>> ['乔治', '佩奇']
\end{lstlisting}

    \item 元组
\begin{lstlisting}[style=demostyle,numbers=none]
变量名 = (值1, 值2, 值3,...)
变量名[起始位置:终止位置:步长]

name = ('张三', '李四', '王五', '乔治', '佩奇')
print(name[1:3])
>> ('李四', '王五')
print(name[3:])
>> ('乔治', '佩奇')
\end{lstlisting}

    \item 字典
\begin{lstlisting}[style=demostyle,numbers=none]
变量名 = {键1:值1, 键2:值2, 键3:值3,...}
变量名[起始位置:终止位置:步长]

name = {'name':'张三', 'age':18, 'sex':'男'}
print(name['name'])
>> 张三
\end{lstlisting}

    \item 集合

    因为集合是无序的，所以不能使用切片。

\end{enumerate}

\subsubsection{split 分割}
\begin{enumerate}[label=\arabic*)]
    \item 字符串

\begin{lstlisting}[style=demostyle,numbers=none]
>>> chinese_zodiac = "鼠牛虎兔龙蛇马羊猴鸡狗猪"
>>> print(chinese_zodiac.split(" "))    # 默认分割符为空格
['鼠', '牛', '虎', '兔', '龙', '蛇', '马', '羊', '猴', '鸡', '狗', '猪']
>>> print(chinese_zodiac.split(" ", 3)) # 最多分割3次
['鼠', '牛', '虎', '兔龙蛇马羊猴鸡狗猪']
>>> english_zodiac = "rat ox tiger rabbit dragon snake horse sheep monkey rooster dog pig"
>>> print(english_zodiac.split(" ", 3))
['rat', 'ox', 'tiger', 'rabbit dragon snake horse sheep monkey rooster dog pig']
\end{lstlisting}

    \item 列表、元组、字典、集合 \textcolor{red}{\textbf{都}}不支持 \codeinline{split} 方法。
\end{enumerate}

\subsubsection{in/not in 成员关系}
\begin{enumerate}[label=\arabic*)]
    \item \codeinline{in}: 判断某个元素是否在序列中
    \item \codeinline{not in}: 判断某个元素是否不在序列中
\end{enumerate}

\begin{enumerate}[label=\arabic*)]
    \item 字符串
\begin{lstlisting}[style=demostyle,numbers=none]
>>> chinese_zodiac = "鼠牛虎兔龙蛇马羊猴鸡狗猪"
>>> print("狗" in chinese_zodiac)
True
>>> print("狗" not in chinese_zodiac)
False
\end{lstlisting}

    \item 列表
\begin{lstlisting}[style=demostyle,numbers=none]
name = ['张三', '李四', '王五']
print('李四' in name)
>> True
print('李四' not in name)
>> False
\end{lstlisting}

    \item 元组
\begin{lstlisting}[style=demostyle,numbers=none]
name = ('张三', '李四', '王五')
print('李四' in name)
>> True
print('李四' not in name)
>> False
\end{lstlisting}

    \item 字典
\begin{lstlisting}[style=demostyle,numbers=none]
name = {'name': '张三', 'age': 18}
print('name' in name)
>> True
print('name' not in name)
>> False
\end{lstlisting}

    \item 集合
\begin{lstlisting}[style=demostyle,numbers=none]
name = {'张三', '李四'}
print('张三' in name)
>> True
print('张三' not in name)
>> False
\end{lstlisting}

\end{enumerate}



\subsubsection{拼接}
\begin{enumerate}[label=\arabic*)]
    \item \codeinline{+}: 将两个序列拼接
    \item \codeinline{join}: 将序列中的元素以指定字符连接
\end{enumerate}
\begin{enumerate}[label=\arabic*)]
    \item 字符串
\begin{lstlisting}[style=demostyle,numbers=none]
>>> a = "123"
>>> b = "456"
>>> print(a + b)
123456
\end{lstlisting}

    \item 列表
\begin{lstlisting}[style=demostyle,numbers=none]
list1 = [1, 2, 3]
list2 = [4, 5, 6]
print(list1 + list2)
>>> [1, 2, 3, 4, 5, 6] 
\end{lstlisting}

    \item 元组
\begin{lstlisting}[style=demostyle,numbers=none]
tuple1 = (1, 2, 3)
tuple2 = (4, 5, 6)
print(tuple1 + tuple2)
>>> (1, 2, 3, 4, 5, 6)
\end{lstlisting}

    \item 字典

    字典不能使用 \codeinline{+} 进行拼接。

    但是字典有\codeinline{|}并集，可以进行拼接。
\begin{lstlisting}[style=demostyle,numbers=none]
set1 = {1, 2, 3}
set2 = {4, 5, 6}
print(set1 | set2)
>>> {1, 2, 3, 4, 5, 6}
\end{lstlisting}

    \item 集合

    集合不能使用 \codeinline{+} 进行拼接。
\end{enumerate}


\subsubsection{join 连接}
\begin{enumerate}[label=\arabic*)]
    \item 字符串
\begin{lstlisting}[style=demostyle,numbers=none]
my_list = ['a', 'b', 'c']
result = ''.join(my_list)  # 使用空字符串连接
print(result)
>>> abc

result_with_dash = '-'.join(my_list)
print(result_with_dash) 
>>> a-b-c
\end{lstlisting}

    \item 列表
\begin{lstlisting}[style=demostyle,numbers=none]
>>> chinese_zodiac = "鼠牛虎兔龙蛇马羊猴鸡狗猪"
>>> print(" ".join(chinese_zodiac))
鼠 牛 虎 兔 龙 蛇 马 羊 猴 鸡 狗 猪
\end{lstlisting}

    \item 元组
\begin{lstlisting}[style=demostyle,numbers=none]
>>> my_tuple = ('a', 'b', 'c')
>>> print(my_tuple * 3)
('a', 'b', 'c', 'a', 'b', 'c', 'a', 'b', 'c')
\end{lstlisting}

    \item 字典
\begin{lstlisting}[style=demostyle,numbers=none]
# 字典
my_dict = {'a': 1, 'b': 2, 'c': 3}
 
# 连接字典的键
keys_result = ''.join(my_dict.keys())
print(keys_result)  
>>> abc
 
# 连接字典的值（需要先将值转换为字符串）
values_result = ''.join(map(str, my_dict.values()))
print(values_result)  
>>> 123
\end{lstlisting}

    \item 集合

\begin{lstlisting}[style=demostyle,numbers=none]
my_set = {1, 2, 3, 4, 5}
 
# 连接集合的元素
elements_result = ''.join(map(str, my_set))
print(elements_result)  
>>> 12345
\end{lstlisting}
\end{enumerate}




\subsubsection{* 重复}
\begin{enumerate}[label=\arabic*)]
    \item 字符串
\begin{lstlisting}[style=demostyle,numbers=none]
>>> chinese_zodiac = "鼠牛虎兔龙蛇马羊猴鸡狗猪"
>>> print(chinese_zodiac * 3)
鼠牛虎兔龙蛇马羊猴鸡狗猪鼠牛虎兔龙蛇马羊猴鸡狗猪鼠牛虎兔龙蛇马羊猴鸡狗猪
\end{lstlisting}

    \item 列表
\begin{lstlisting}[style=demostyle,numbers=none]
my_list = [1, 2, 3]
print(my_list * 2)  
>>> [1, 2, 3, 1, 2, 3]
\end{lstlisting}

    \item 元组
\begin{lstlisting}[style=demostyle,numbers=none]
my_tuple = (1, 2, 3)
print(my_tuple * 2)  
>>> (1, 2, 3, 1, 2, 3)
\end{lstlisting}

    \item 字典

    字典不支持重复操作。

    \item 集合

    集合不支持重复操作。

\end{enumerate}

\subsubsection{append 追加}
\begin{enumerate}[label=\arabic*)]
    \item 字符串

    字符串不支持 \codeinline{append} 方法。

    \item 列表
\begin{lstlisting}[style=demostyle,numbers=none]
my_list = [1, 2, 3]
my_list.append(4)
print(my_list)  
>>> [1, 2, 3, 4]
\end{lstlisting}

    \item 元组

    元组不支持 \codeinline{append} 方法。

    \item 字典
\begin{lstlisting}[style=demostyle,numbers=none]
my_dict = {'a': 1, 'b': 2}
my_dict['c'] = 3
print(my_dict)  
>>> {'a': 1, 'b': 2, 'c': 3}
\end{lstlisting}

    \item 集合
\begin{lstlisting}[style=demostyle,numbers=none]
my_set = {1, 2, 3}
my_set.add(4)
print(my_set)  
>>> {1, 2, 3, 4}
\end{lstlisting}
\end{enumerate}

\subsubsection{extend 扩展}
\begin{enumerate}[label=\arabic*)]
    \item 字符串

\begin{lstlisting}[style=demostyle,numbers=none]
my_str = 'abc'
my_str.extend('def')
print(my_str)  
>>> abcdef
\end{lstlisting}

    \item 列表

    \codeinline{extend}方法 \textcolor{red}{\textbf{只}}接收列表中的元素逐一添加到另一个列表中。

    该方法会将参数列表中的元素逐一添加到原列表中，而不是将参数列表作为一个整体添加到原列表中。
\begin{lstlisting}[style=demostyle,numbers=none]
my_list = [1, 2, 3]
my_list.extend([4, 5])
print(my_list)  
>>> [1, 2, 3, 4, 5]
\end{lstlisting}

    \item 元组

    元组不支持 \codeinline{extend} 方法。

    \item 字典
\begin{lstlisting}[style=demostyle,numbers=none]
my_dict = {'a': 1, 'b': 2}
my_dict.update({'c': 3, 'd': 4})
print(my_dict)  
>>> {'a': 1, 'b': 2, 'c': 3, 'd': 4}
\end{lstlisting}

    \item 集合
\begin{lstlisting}[style=demostyle,numbers=none]
my_set = {1, 2, 3}
my_set.update({4, 5})
print(my_set)  
>>> {1, 2, 3, 4, 5}
\end{lstlisting}
\end{enumerate}

\subsubsection{insert 插入}
\begin{enumerate}[label=\arabic*)]
    \item 字符串

    字符串不支持 \codeinline{insert} 方法。

    \item 列表

\begin{lstlisting}[style=demostyle,numbers=none]
my_list = [1, 2, 3]
my_list.insert(1, 1.5)
print(my_list)  
>>> [1, 1.5, 2, 3]
\end{lstlisting}

    \item 元组

    元组不支持 \codeinline{insert} 方法。

    \item 字典

\begin{lstlisting}[style=demostyle,numbers=none]
my_dict = {'a': 1, 'b': 2}
my_dict.insert('c', 3)
print(my_dict)  
>>> {'a': 1, 'b': 2}
\end{lstlisting}

    \item 集合

    集合不支持 \codeinline{insert} 方法。
\end{enumerate}

\subsubsection{update 更新}
\begin{enumerate}[label=\arabic*)]
    \item 字符串

    字符串是不可变的，因此没有 \codeinline{update} 方法。

    \item 列表

    列表的 \codeinline{update} 方法是 \codeinline{extend} 方法。

    \item 元组

    元组不支持 \codeinline{update} 方法。

    \item 字典
\begin{lstlisting}[style=demostyle,numbers=none]
my_dict = {'a': 1, 'b': 2}
my_dict.update({'b': 3, 'c': 4})  # 更新键 'b' 的值，并添加键 'c'
print(my_dict) 
>>> {'a': 1, 'b': 3, 'c': 4}
\end{lstlisting}

    \item 集合

\begin{lstlisting}[style=demostyle,numbers=none]
my_set = {1, 2, 3}
my_set.update({3, 4, 5})  # 添加元素 4 和 5
print(my_set)  
>>> {1, 2, 3, 4, 5}
\end{lstlisting}
\end{enumerate}

\subsubsection{remove 删除}
\begin{enumerate}[label=\arabic*)]
    \item 字符串

    字符串不支持 \codeinline{remove} 方法。

    \item 列表

\begin{lstlisting}[style=demostyle,numbers=none]
my_list = [1, 2, 3, 2]
my_list.remove(2)
print(my_list)  
>>> [1, 3, 2]
\end{lstlisting}

    \item 元组

    元组不支持 \codeinline{remove} 方法。

    \item 字典
\begin{lstlisting}[style=demostyle,numbers=none]
my_dict = {'a': 1, 'b': 2}
del my_dict['b']
print(my_dict)
>>> {'a': 1}
\end{lstlisting}

    \item 集合
\begin{lstlisting}[style=demostyle,numbers=none]
my_set = {1, 2, 3, 2}
my_set.remove(2)  # 删除第一个匹配的2
print(my_set)  # 输出: {1, 3, 2}
\end{lstlisting}

    集合还有一个\codeinline{discard}方法，如果元素不存在，\codeinline{remove}会报错，而\codeinline{discard}不会。
\begin{lstlisting}[style=demostyle,numbers=none]
my_set = {1, 2, 3, 4, 5}
my_set.discard(4)
print(my_set)
>>> {1, 2, 3, 5}
\end{lstlisting}
\end{enumerate}

\subsubsection{pop 弹出}
\begin{enumerate}[label=\arabic*)]
    \item 字符串

    字符串不支持 \codeinline{pop} 方法。

    \item 列表

\begin{lstlisting}[style=demostyle,numbers=none]
my_list = [1, 2, 3, 2]
my_list.remove(2)  # 删除第一个匹配的2
print(my_list)  # 输出: [1, 3, 2]
my_list.pop(1)  # 删除索引为1的元素
print(my_list)  # 输出: [1, 2]
\end{lstlisting}

    \item 元组

    元组不支持 \codeinline{pop} 方法。

    \item 字典

\begin{lstlisting}[style=demostyle,numbers=none]
my_dict = {'a': 1, 'b': 2}
my_dict.pop('a')
>>> 1
\end{lstlisting}

    \item 集合

    集合的 \codeinline{pop} 方法是随机删除一个元素。
\begin{lstlisting}[style=demostyle,numbers=none]
## pop方法:随机删除一个元素
set1 = {1, 2, 3, 4, 5, 6, 7}
set1.pop()
print(set1)
>>> {2, 3, 4, 5, 6, 7}
\end{lstlisting}
\end{enumerate}



\subsubsection{clear 清空}
\begin{enumerate}[label=\arabic*)]
    \item 字符串

    字符串是不可变的，因此没有 \codeinline{clear} 方法。

    \item 列表

\begin{lstlisting}[style=demostyle,numbers=none]
## clear方法:清空列表
list1 = [1, 2, 3, 4, 5, 6, 7]
list1.clear()
print(list1)
>>> []
\end{lstlisting}

    \item 元组

    元组是不可变的，因此没有 \codeinline{clear} 方法。

    \item 字典
\begin{lstlisting}[style=demostyle,numbers=none]
## clear方法:清空字典
dict1 = {'name': 'Jack', 'age': 20, 'gender': 'male'}
dict1.clear()
print(dict1)
>>> {}
\end{lstlisting}

    \item 集合

\begin{lstlisting}[style=demostyle,numbers=none]
my_set = {1, 2, 3, 4, 5}
my_set.clear()  # 清空集合
print(my_set)  
>>> set()
\end{lstlisting}
\end{enumerate}

\subsubsection{del 删除}
实际开发中并不经常使用 \codeinline{del} 方法，因为 Python 自带的垃圾回收机制会自动销毁无用的对象。

\begin{enumerate}[label=\arabic*)]
    \item 字符串

    字符串是不可变的，因此没有 \codeinline{del} 方法。

    \item 列表
\begin{lstlisting}[style=demostyle,numbers=none]
intlist = [1, 45, 8, 34]
print(intlist)
del intlist[1]
print(intlist)
>>> [1, 8, 34]
del intlist
print(intlist)
>>> NameError: name 'intlist' is not defined
\end{lstlisting}

    \item 元组

\begin{lstlisting}[style=demostyle,numbers=none]
my_tuple = (1, 2, 3, 4, 5)
print(my_tuple)
>>> (1, 2, 3, 4, 5)
del my_tuple
print(my_tuple)
>>> NameError: name 'my_tuple' is not defined
\end{lstlisting}

    \item 字典

\begin{lstlisting}[style=demostyle,numbers=none]
my_dict = {'a': 1, 'b': 2, 'c': 3}
del my_dict['b']  # 删除键为 'b' 的键值对
print(my_dict)
>>> {'a': 1, 'c': 3}

del my_dict
print(my_dict)
>>> NameError: name 'my_dict' is not defined
\end{lstlisting}

    \item 集合

\begin{lstlisting}[style=demostyle,numbers=none]
my_set = {1, 2, 3, 4, 5}
del my_set
print(my_set)
>>> NameError: name 'my_set' is not defined
\end{lstlisting}

\end{enumerate}



\subsection{序列间转换}
\subsubsection{list 将序列转为列表}
\begin{enumerate}[label=\arabic*)]
    \item 字符串
\begin{lstlisting}[style=demostyle,numbers=none]
my_string = "hello"
print(list(my_string))  # 输出: ['h', 'e', 'l', 'l', 'o']
>>> ['h', 'e', 'l', 'l', 'o']
\end{lstlisting}

    \item 列表

    列表不存在转换为列表的问题，所以转换后的结果还是原列表。

    \item 元组
\begin{lstlisting}[style=demostyle,numbers=none]
my_tuple = (1, 2, 3)
print(list(my_tuple))  # 输出: [1, 2, 3]
>>> [1, 2, 3]
\end{lstlisting}

    \item 字典
\begin{lstlisting}[style=demostyle,numbers=none]
my_dict = {'a': 1, 'b': 2, 'c': 3}
print(list(my_dict))  # 输出: ['a', 'b', 'c']
>>> ['a', 'b', 'c']
\end{lstlisting}

    \item 集合
\begin{lstlisting}[style=demostyle,numbers=none]
my_set = {1, 2, 3}
print(list(my_set))  # 输出: [1, 2, 3]
>>> [1, 2, 3]
\end{lstlisting}
\end{enumerate}

\subsubsection{tuple 将序列转为元组}
\begin{enumerate}[label=\arabic*)]
    \item 字符串
\begin{lstlisting}[style=demostyle,numbers=none]
my_string = "hello"
print(tuple(my_string)) 
>>> ('h', 'e', 'l', 'l', 'o')
\end{lstlisting}

    \item 列表
\begin{lstlisting}[style=demostyle,numbers=none]
my_list = [1, 2, 3]
print(tuple(my_list))
>>> (1, 2, 3)
\end{lstlisting}

    \item 元组

    元组不存在转换为元组的问题，所以转换后的结果还是原元组。

    \item 字典
\begin{lstlisting}[style=demostyle,numbers=none]
my_dict = {'a': 1, 'b': 2}
print(tuple(my_dict))
>>> ('a', 'b')
\end{lstlisting}

    \item 集合
\begin{lstlisting}[style=demostyle,numbers=none]
my_set = {1, 2, 3}
print(tuple(my_set))
>>> (1, 2, 3)
\end{lstlisting}

\end{enumerate}

\subsubsection{set 将序列转为集合}
\begin{enumerate}[label=\arabic*)]
    \item 字符串
\begin{lstlisting}[style=demostyle,numbers=none]
my_string = "hello"
print(set(my_string))
>>> {'h', 'e', 'l', 'o'}
\end{lstlisting}

    \item 列表
\begin{lstlisting}[style=demostyle,numbers=none]
my_list = [1, 2, 3]
print(set(my_list))
>>> {1, 2, 3}
\end{lstlisting}

    \item 元组
\begin{lstlisting}[style=demostyle,numbers=none]
my_tuple = (1, 2, 3)
print(set(my_tuple))
>>> {1, 2, 3}
\end{lstlisting}
    
    \item 字典
\begin{lstlisting}[style=demostyle,numbers=none]
my_dict = {'a': 1, 'b': 2}
print(set(my_dict))
>>> {'b', 'a'}
\end{lstlisting}

    \item 集合

    集合不存在转换为集合的问题，所以转换后的结果还是原集合。

\end{enumerate}





\section{列表(list)}
\subsection{创建}
\subsubsection{直接生成}
\begin{lstlisting}[style=demostyle,numbers=none]
变量名 = [值1, 值2, 值3,...]
\end{lstlisting}

\begin{lstlisting}[style=demostyle,numbers=none]
name = ['张三', '李四', '王五', 1, 2.3, True]
>>> ['张三', '李四', '王五', 1, 2.3, True]
\end{lstlisting}

\subsubsection{list 显式转换}
\codeinline{range} 函数用于生成一个数字序列，可以通过 \codeinline{list} 函数将其转换为列表。
\begin{lstlisting}[style=demostyle,numbers=none]
变量名 = list(range(起始值, 终止值, 步长))
\end{lstlisting}

\begin{lstlisting}[style=demostyle,numbers=none]
>>> list(range(10))
[0, 1, 2, 3, 4, 5, 6, 7, 8, 9]
>>> list(range(1, 10))
[1, 2, 3, 4, 5, 6, 7, 8, 9]
>>> list(range(1, 10, 2))
[1, 3, 5, 7, 9]
\end{lstlisting}

\subsubsection{列表推导式}
\begin{lstlisting}[style=demostyle,numbers=none]
变量名 = [表达式 for 变量 in 可迭代对象 (if 条件)]
\end{lstlisting}

\begin{lstlisting}[style=pythonstyle]
a = [1, 2, 3, 4, 5]
b = [i * 2 for i in a]
print(b)
>> [2, 4, 6, 8, 10]
\end{lstlisting}

\begin{lstlisting}[style=pythonstyle]
# 生成从1-10所有偶数的平方
alist = [i**2 for i in range(1, 11) if i % 2 == 0]
print(alist)
>> [4, 16, 36, 64, 100]
\end{lstlisting}

\begin{lstlisting}[style=pythonstyle]
# 过滤掉长度小于或等于3的字符串列表，并将剩下的转换成大写字母：
names = ['Bob','Tom','alice','Jerry','Wendy','Smith']
new_names = [name.upper() for name in names if len(name)>3]
print(new_names)
>>> ['ALICE', 'JERRY', 'WENDY', 'SMITH']
\end{lstlisting}

\begin{lstlisting}[style=pythonstyle]
# 计算 30 以内可以被 3 整除的整数：
multiples = [i for i in range(30) if i % 3 == 0]
print(multiples)
>>> [0, 3, 6, 9, 12, 15, 18, 21, 24, 27]
\end{lstlisting}

\subsection{修改}
直接对元素赋值即可
\begin{lstlisting}[style=demostyle,numbers=none]
a = [1, 2, 3]
a[0] = 4
print(a)
>> [4, 2, 3]
\end{lstlisting}

\subsection{栈}
使用 list 列表模拟栈功能的实现方法是，
\begin{enumerate}[label=\arabic*)]
    \item 使用 \codeinline{append()} 方法\textcolor{YBXPurple}{在末尾}存入数据；
    \item 使用 \codeinline{pop()} 方法\textcolor{YBXPurple}{在末尾}读取数据。
\end{enumerate}

\begin{lstlisting}[style=demostyle,numbers=none]
stack = []
stack.append(1)
stack.append("hello")
stack.append(2)
stack.append(3)
print(stack)
>> [1, 'hello', 2, 3]

print(stack.pop())
>> 3

print(stack.pop())
>> 2

print(stack.pop())
>> hello
\end{lstlisting}

\subsection{collections 实现栈和队列}
面使用 list 实现队列的例子中，插入数据的部分是通过 insert() 方法实现的，这种方法效率并不高，因为每
次从列表的开头插入一个数据，列表中所有元素都得向后移动一个位置。

标准库的 \codeinline{collections} 模块中的 \codeinline{deque} 结构体，它被设计成在两端存
入和读取都很快的特殊 list，可以用来实现栈和队列的功能。



\section{字符串}
\subsection{定义}
字符串是 Python 中最常用的数据类型。我们可以使用引号('或")来创建字符串。\textcolor{YBXPurple}{不能修改}。

\section{集合}
\subsection{定义}
集合是一个 \textcolor{YBXPurple}{无序的、不重复}的数据集合。

\subsubsection{集合推导式}
\begin{lstlisting}[style=demostyle,numbers=none]
变量名 = {表达式 for 变量 in 可迭代对象 (if 条件)}
\end{lstlisting}

\begin{lstlisting}[style=pythonstyle]
# 计算数字 1,2,3 的平方数
num_set = {i**2 for i in (1,2,3)}
num_set
>>> {1, 4, 9}
\end{lstlisting}

\begin{lstlisting}[style=pythonstyle]
# 判断不是 abc 的字母并输出：
letters = 'abcdefghijklmnopqrstuvwxyz'
not_abc = {x for x in letters if x not in 'abc'}
print(not_abc)
>>> {'u', 'j', 'p', 'i', 'r', 'f', 'q', 'v', 'x', 'y', 'z', 'h', 'l', 't', 'o', 'w', 'm', 'e', 'n', 'g', 'k', 's', 'd'}
\end{lstlisting}

\subsection{集合间运算}
\textcolor{red}{\textbf{不同于}}其他容器类型(如列表、元组等)，集合类型支持集合间的运算。

\begin{lstlisting}[style=demostyle,numbers=none]
a = {1, 2, 3, 4, 5}
b = {4, 5, 6, 7, 8}

# 交集
## & 符号
print(a & b)
>>> {4, 5}
## intersection方法
print(a.intersection(b))
>>> {4, 5}

# 并集
## | 符号
print(a | b)
>>> {1, 2, 3, 4, 5, 6, 7, 8}
## union方法
print(a.union(b))
>>> {1, 2, 3, 4, 5, 6, 7, 8}

# 差集
## - 符号
print(a - b)
>>> {1, 2, 3}
print(b - a)
>>> {8, 6, 7}
## difference方法
print(a.difference(b))
>>> {1, 2, 3}
print(b.difference(a))
>>> {8, 6, 7}

# 对称差集
## ^ 符号
print(a ^ b)
>>> {1, 2, 3, 6, 7, 8}
## symmetric_difference方法
print(a.symmetric_difference(b))
>>> {1, 2, 3, 6, 7, 8}
\end{lstlisting}

\section{字典}
字典是一种无序的数据集合，使用键-值（key-value）存储，具有极快的查找速度。
\begin{lstlisting}[style=demostyle,numbers=none]
# 字典数据没有下标，通过键来查找

变量名 = {键1:值1, 键2:值2, 键3:值3,...}
\end{lstlisting}


\subsection{创建}
\subsubsection{直接生成}
\begin{lstlisting}[style=demostyle,numbers=none]
# 增加
dict1 = {'name': '张三', 'age': 18}

# 添加键值对
dict1['sex'] = '男'
print(dict1)
>>> {'name': '张三', 'age': 18, 'sex': '男'}

# 修改
dict1['age'] = 20
print(dict1)
>>> {'name': '张三', 'age': 20, 'sex': '男'}
\end{lstlisting}


\subsubsection{字典推导式}
\begin{lstlisting}[style=demostyle,numbers=none]
变量名 = {键:值 for 变量 in 可迭代对象 (if 条件)}
\end{lstlisting}

\begin{lstlisting}[style=pythonstyle]
a = [1, 2, 3, 4, 5]
b = {i: i * 2 for i in a}
print(b)
>> {1: 2, 2: 4, 3: 6, 4: 8, 5: 10}
\end{lstlisting}

\begin{lstlisting}[style=pythonstyle]
# 生成从1-10所有偶数的平方
adict = {i: i**2 for i in range(1, 11) if i % 2 == 0}
print(adict)
>>> {2: 4, 4: 16, 6: 36, 8: 64, 10: 100}
\end{lstlisting}

\begin{lstlisting}[style=pythonstyle]
names = ['Bob','Tom','alice','Jerry','Wendy','Smith']    
# 将列表中各字符串值为键，各字符串的长度为值，组成键值对
newdict = {key:len(key) for key in names}
print(newdict)
>>> 
{'Bob': 3, 'Tom': 3, 'alice': 5, 'Jerry': 5, 'Wendy': 5, 'Smith': 5}
\end{lstlisting}


\section{元组(tuple)}
\subsection{定义}
元组与列表类似，不同之处在于元组的元素 \textcolor{YBXPurple}{不能修改}。

\subsection{创建}
\subsubsection{直接生成}
\begin{lstlisting}[style=demostyle,numbers=none]
# 直接创建元组
t2 = (1, 2, 3)
print(t2)
>>> (1, 2, 3)
\end{lstlisting}

\subsubsection{tuple 显式转换}
\begin{lstlisting}[style=demostyle,numbers=none]
# 通过 tuple 方法创建元组
t1 = tuple([1, 2, 3])
print(t1)
>>> (1, 2, 3)
\end{lstlisting}

\subsubsection{元组推导式(生成器表达式)}
\begin{lstlisting}[style=demostyle,numbers=none]
变量名 = (表达式 for 变量 in 可迭代对象 (if 条件))
\end{lstlisting}

\begin{lstlisting}[style=pythonstyle]
a = (x for x in range(1,10))
print(a)
>>> <generator object <genexpr> at 0x7f16c56bf840>
\end{lstlisting}
\textcolor{red}{\textbf{不同于列表与字典，元组推到其的结果是生成器}}，可以使用\codeinline{tuple()}方法将生成器转换为元组。
\begin{lstlisting}[style=pythonstyle]
a = (x for x in range(1,10))
a = tuple(a)
print(a)
>>> (1, 2, 3, 4, 5, 6, 7, 8, 9)
\end{lstlisting}